Dans le référentiel $S$ d'un laboratoire, l'événement A se produit à $x_A = \qty{0}{m}$, $t_A = \qty{0}{s}$ et l'événement B se produit à $x_B = \qty{\v{xB}}{m}$, $t_B = \qty{\v{tB}}{\micro\second}$. Un référentiel $S'$ se déplace à $\num{\v{beta}}\,c$ selon l'axe des $x$ positifs par rapport à $S$ ; les origines coïncident à $t = t' = 0$.

\begin{enumerate}[label=\alph*)]
    \item Déterminez le facteur de Lorentz.
    \item À l'aide des transformations de Lorentz, déterminez la position $x'_B$ et l'instant $t'_B$ de l'événement B dans $S'$.
    \item Vérifiez que la quantité $(\Delta x)^2 - (c\,\Delta t)^2$ a la même valeur dans les deux référentiels. Les deux événements peuvent-ils avoir un lien de cause à effet ?
\end{enumerate}
